\chapter*{第一部分学习路线}
\addcontentsline{toc}{chapter}{第一部分学习路线：从初中语言进入高中模型}
\markboth{第一部分学习路线}{第一部分学习路线}

\begin{learningbox}{从初中语言进入高中模型}
这一部分的重点不是多做难题，而是完成一次语言转换：从初中常说的“快慢、远近、受力、运动状态”，转到高中要使用的“位移、速度、加速度、合力、模型、图像和正负号”。
\end{learningbox}

\section*{一、这一部分的认知台阶}

\begin{enumerate}[itemsep=0.2em]
  \item \textbf{先学物理语言：} Day 01 建立模型、矢量、单位和量纲意识。
  \item \textbf{再学运动描述：} Day 02 到 Day 04 用位移、速度、加速度和图像描述运动。
  \item \textbf{最后进入受力分析：} Day 05 到 Day 07 把力、牛顿定律和连接体合在一起。
\end{enumerate}

这三个台阶不能颠倒。若还分不清位移和路程、速度和速率，后面的加速度、图像和牛顿第二定律就容易混乱。

\section*{二、阅读时抓住三条线}

\begin{itemize}[itemsep=0.2em]
  \item \textbf{方向线：} 位移、速度、加速度和力都是矢量。先规定正方向，再把方向写进正负号。
  \item \textbf{图像线：} $x$-$t$ 图看斜率，$v$-$t$ 图看斜率和面积。
  \item \textbf{模型线：} 质点、光滑面、轻绳、轻滑轮、整体法和隔离法，都是把复杂真实情境简化成可计算问题。
\end{itemize}

\section*{三、卡住时先回到哪里}

运动题卡住，先回 Day 02 看清坐标轴、位移和速度；匀变速与图像题卡住，回 Day 03、Day 04 找正方向、斜率和面积；受力题卡住，回 Day 05、Day 06 先画受力图。
